\documentclass[11pt,letterpaper]{article}
\usepackage[margin=1in]{geometry}
\usepackage{amsmath,amssymb,amsthm}
\usepackage{booktabs}
\usepackage{enumitem}
\usepackage{hyperref}
\usepackage{xcolor}

\hypersetup{colorlinks=true,linkcolor=blue!60!black,citecolor=blue!60!black,urlcolor=blue!60!black}

\newtheorem{remark}{Remark}

\title{The DSAT Recreational Dive Planner:\\
       Algorithm, Mathematics, and Physical Justification\\[6pt]
       \large Rebuilt from Hamilton, Rogers, Powell, \& Vann (1994)}
\author{Christopher}
\date{March 2026}

\begin{document}
\maketitle

\begin{abstract}
This document presents a complete, self-contained derivation of the decompression
algorithm underlying PADI's Recreational Dive Planner (RDP), reconstructed from
the primary source report by Hamilton, Rogers, Powell, and Vann (1994).
Every parameter, equation, and design decision is stated explicitly, traced to
its physical or empirical justification, and verified numerically against the
published tables.  The goal is an independent audit: a diver should be able
to read this document, understand every step, and reproduce the published
NDL tables to within the known fitting residuals.
\end{abstract}

\tableofcontents
\newpage

% ==================================================================
\section{Physical Foundations}
% ==================================================================

\subsection{The dissolved-gas problem}

A scuba diver breathes compressed air at ambient pressure.  The partial
pressure of nitrogen in the inspired gas increases with depth, driving
nitrogen into solution in the blood and tissues.  On ascent, the ambient
pressure drops and the dissolved nitrogen becomes supersaturated.  If
the supersaturation is too large, nitrogen comes out of solution as
bubbles, potentially causing decompression sickness (DCS).

The engineering problem is: given a time-depth profile, predict whether
the dissolved nitrogen in any body tissue will exceed a safe
supersaturation limit at any point during or after ascent.

\subsection{Units and conventions}

All pressures are expressed in \textbf{feet of seawater} (fsw), defined as
\[
  1\;\text{fsw} = \frac{1}{33}\;\text{atm} = 3.070\;\text{kPa}.
\]
One atmosphere of surface pressure equals $33$~fsw.  The metric unit
\textbf{msw} (metres of seawater) is defined as $1\;\text{msw} = \tfrac{1}{10}\;\text{bar} = 10\;\text{kPa}$.
The conversion is
$1\;\text{msw} = 3.2568\;\text{fsw}$
(Hamilton et al.\ 1994, p.~8).

\textbf{Absolute pressure} at gauge depth $d$ (fsw):
\begin{equation}
  P_{\mathrm{amb}} = d + 33 \quad (\text{fsw}).
  \label{eq:pamb}
\end{equation}

\subsection{Gas composition}

The DSAT model treats air as $79.1\%$ nitrogen by volume.  The inspired
nitrogen partial pressure at ambient pressure $P_{\mathrm{amb}}$ is
\begin{equation}
  P_i = f_{\mathrm{N}_2} \cdot P_{\mathrm{amb}},
  \qquad f_{\mathrm{N}_2} = 0.791.
  \label{eq:pi}
\end{equation}

\begin{remark}[No water-vapor correction]
Unlike the B\"uhlmann ZH-L16 model, the DSAT algorithm does \emph{not}
subtract alveolar water-vapor pressure before computing $P_i$.  The report
states: ``exchange is between the compartments and the ambient atmosphere
\ldots that factor is not used'' (Hamilton et al., p.~13).  This is
self-consistent because the $M_0$ values were back-calculated under the
same convention.
\end{remark}

\begin{remark}[Historical note on the nitrogen fraction]
Haldane (1908) and the original USN tables treated air as $100\%$ nitrogen
to simplify hand calculation.  The $M$-values were scaled accordingly, so
the ratio-based ascent limits gave the same answer.  Rogers adopted the
modern convention of using the true nitrogen fraction $0.791$ with
correspondingly lower $M$-values (Hamilton et al., pp.~13--14, 21--22).
\end{remark}


% ==================================================================
\section{The Haldane Compartment Model}
% ==================================================================

\subsection{Parallel perfusion-limited compartments}

The body is modeled as $N=14$ independent, parallel compartments, each
exchanging nitrogen with the lungs at a rate determined by a single
\textbf{half-time} $t_{1/2}$.  A compartment is not an anatomical structure;
it represents all tissues in the body that happen to equilibrate on a
similar timescale.

\subsection{Gas loading at constant depth (Haldane equation)}

Let $P_0$ be the nitrogen partial pressure in a compartment at time $t=0$,
and $P_i$ the constant inspired nitrogen pressure.  The compartment
loading at time $t$ is
\begin{equation}
  \boxed{
  P(t) = P_0 + (P_i - P_0)\bigl(1 - e^{-kt}\bigr)
  }
  \label{eq:haldane}
\end{equation}
where the time constant is
\begin{equation}
  k = \frac{\ln 2}{t_{1/2}}.
  \label{eq:k}
\end{equation}

\textbf{Physical justification.}  Equation~\eqref{eq:haldane} is the
solution to the first-order ODE
\[
  \frac{dP}{dt} = k\,(P_i - P),
\]
which states that the rate of gas exchange is proportional to the partial
pressure gradient between the inspired gas and the tissue.  This is the
defining assumption of a \emph{perfusion-limited} model: blood flow
(perfusion) is the rate-limiting step, not diffusion within tissue.  The
exponential approach to equilibrium is a direct consequence.

\textbf{Key property:} After one half-time, $50\%$ of the pressure difference
has been absorbed.  After six half-times, $98.4\%$---the compartment is
considered saturated.

\subsection{Gas loading during linear pressure change (Schreiner equation)}

During descent or ascent at a constant rate, the ambient pressure changes
linearly and so does $P_i(t) = P_i(0) + Rt$, where $R = f_{\mathrm{N}_2}
\times (\text{depth rate in fsw/min})$.  The Haldane ODE becomes
\[
  \frac{dP}{dt} = k\bigl(P_i(0) + Rt - P\bigr),
\]
which has the closed-form solution (Schreiner and Kelley, 1971):
\begin{equation}
  \boxed{
  P(t) = P_i(0) + R\!\left(t - \frac{1}{k}\right)
       - \left[P_i(0) - P_0 - \frac{R}{k}\right] e^{-kt}
  }
  \label{eq:schreiner}
\end{equation}
This is used to model finite-rate descent ($R > 0$) and ascent ($R < 0$).

\subsection{Compartment parameters}

The 14 DSAT half-times and their corresponding surfacing $M$-values
($M_0$, in fsw) are given in Table~\ref{tab:params}.  These are taken
directly from Table~III of Hamilton et al.\ (1994), ``RDP'' column.

\begin{table}[h]
\centering
\caption{DSAT compartment parameters (Hamilton et al.\ 1994, Table~III).}
\label{tab:params}
\begin{tabular}{@{}rrc@{}}
\toprule
Compartment & $t_{1/2}$ (min) & $M_0$ (fsw) \\
\midrule
 1 &   5 & 99.08 \\
 2 &  10 & 82.63 \\
 3 &  20 & 66.89 \\
 4 &  30 & 59.74 \\
 5 &  40 & 55.73 \\
 6 &  60 & 51.44 \\
 7 &  80 & 49.21 \\
 8 & 100 & 47.85 \\
 9 & 120 & 46.93 \\
10 & 160 & 45.78 \\
11 & 200 & 45.07 \\
12 & 240 & 44.60 \\
13 & 360 & 43.81 \\
14 & 480 & 43.40 \\
\bottomrule
\end{tabular}
\end{table}


% ==================================================================
\section{Ascent Limits: $M_0$ Values}
% ==================================================================

\subsection{The concept of $M$-values}

An $M$-value (Workman, 1965) is the maximum tolerable inert-gas partial
pressure in a compartment at a given ambient pressure.  For no-stop
(no-decompression) diving, only the \textbf{surfacing} $M$-value, $M_0$,
matters: it is the maximum nitrogen loading that a compartment may
carry when the diver surfaces to 1~atm ($33$~fsw ambient).

If any compartment's loading $P_i > M_{0,i}$ at the surface, the model
predicts an unacceptable risk of DCS.  For no-stop planning, the diver
must surface before any compartment reaches its $M_0$.

\subsection{How the DSAT $M_0$ values were derived}
\label{sec:m0derivation}

The DSAT $M_0$ values were \emph{not} chosen from first principles or from
a physiology-based bubble model.  The derivation, stated explicitly in the
report (Hamilton et al., pp.~20--21), was:

\begin{enumerate}[label=(\arabic*)]
  \item \textbf{Empirical NDLs from Doppler data.}
    Spencer (1976) measured Doppler-detectable circulating bubbles in divers
    after no-stop dives at various depths and times.  He proposed that a
    decompression would be acceptable if no more than $\sim\!20\%$ of
    divers produced detectable bubbles.

  \item \textbf{Power-law smoothing.}
    Rogers fitted the Spencer-adjusted NDLs to a smooth depth--time curve
    of the Hempleman form:
    \begin{equation}
      D - A = C\, t^x,
      \label{eq:nostop_curve}
    \end{equation}
    where $D$ is gauge depth (fsw), $t$ is allowable time (min),
    $A = 20.15$~fsw (the shallow asymptote---the shallowest depth at which
    a diver could saturate and still require a stop), $C = 510$, and
    $x = -0.53$ (nearly a square-root law).  This was set slightly
    more conservative than Spencer's suggestion.

  \item \textbf{Back-calculation of $M_0$.}
    For each compartment $i$, the $M_0$ value was determined by asking:
    ``What nitrogen loading would compartment $i$ have if the diver
    descended at 60~fsw/min and remained at the depth where compartment
    $i$ is controlling for exactly the NDL given by the smoothed curve?''
    Solving the Haldane equation forward from surface saturation to the
    NDL time and reading off $P(t_{\mathrm{NDL}})$ gives $M_{0,i}$.
\end{enumerate}

\begin{remark}[The causal arrow]
The $M_0$ values are \emph{derived quantities}, not primitive inputs.
The primary data are the empirical NDLs.  This was confirmed by
Dr.~Michael Powell (co-developer of the RDP) on ScubaBoard (2008):
``One uses the NDLs to calculate a `limiting' halftime for each of
these depths \ldots\ The $M_0$-values will depend on the smoothed NDL
curve.''
\end{remark}


% ==================================================================
\section{No-Decompression Limit Calculation}
% ==================================================================

\subsection{NDL for a single compartment}

Given inspired pressure $P_i$, initial loading $P_0$, and surfacing limit
$M_{0}$, solve $P(t) = M_0$ in~\eqref{eq:haldane} for $t$:
\begin{equation}
  \boxed{
  t_{\mathrm{NDL}} = t_{1/2}\, \log_2\!\left(\frac{P_i - P_0}{P_i - M_0}\right)
  }
  \label{eq:ndl}
\end{equation}
This is valid only when $P_i > M_0$ (the compartment can actually reach
its limit at this depth).  If $P_i \le M_0$, the compartment can never
exceed its limit at this depth, so its NDL is infinite.

\subsection{NDL for the full model}

The overall NDL is the \emph{minimum} across all 14 compartments:
\[
  \mathrm{NDL}(d) = \min_{i=1}^{14}\; t_{\mathrm{NDL},i}(d).
\]
The compartment achieving the minimum is called the
\textbf{controlling compartment}.  As depth increases, the controlling
compartment shifts from slow (large $t_{1/2}$) to fast (small $t_{1/2}$)
compartments.

\subsection{Initial conditions}

For a clean (first dive of the day) diver, all compartments are at
surface-equilibrium nitrogen pressure:
\[
  P_0 = f_{\mathrm{N}_2} \times P_{\text{surface}} = 0.791 \times 33 = 26.10\;\text{fsw}.
\]
For a repetitive dive, $P_0$ for each compartment is the residual loading
from the previous dive(s) after the surface interval.


% ==================================================================
\section{Repetitive Diving}
% ==================================================================

\subsection{Motivation}

After surfacing from a dive, nitrogen continues to off-gas.  A second dive
started before full equilibration begins with elevated tissue loadings,
reducing the available NDL.  The RDP tracks this via pressure groups and
the surface-interval washout model.

\subsection{Pressure groups}

The excess nitrogen in the \textbf{60-minute compartment} (compartment 6,
$t_{1/2} = 60$~min) above surface equilibrium is divided into 26 equal
bins labeled A through Z.  The maximum possible excess is
$M_{0,6} - P_{\text{surface}} \cdot f_{\mathrm{N}_2} = 51.44 - 26.10 = 25.34$~fsw.
Each pressure-group bin spans $25.34/26 \approx 0.975$~fsw.

\textbf{Why the 60-minute compartment?}
The USN tables used the 120-minute compartment for repetitive groups, which
was appropriate for decompression diving.  Rogers argued that for no-stop
recreational dives (which never load the slow compartments to their limits),
the 60-minute compartment better captures the relevant gas dynamics and
provides more efficient surface-interval credit (Hamilton et al., pp.~16--17, 22).

\subsection{Surface washout rule}

During a surface interval, all compartments off-gas toward the surface
equilibrium pressure.  The DSAT model applies a special rule:

\begin{quote}
\emph{At the surface, all compartments with $t_{1/2} \le 60$~min wash out
at the 60-minute rate.  Compartments with $t_{1/2} > 60$~min wash out at
their own (slower) rate.}
\end{quote}

Mathematically, during a surface interval of duration $t_s$:
\begin{equation}
  P_i^{\text{after}} = P_{\text{eq}} + (P_i^{\text{before}} - P_{\text{eq}})
    \cdot 2^{-t_s / t_{1/2}^{\text{eff}}},
  \label{eq:washout}
\end{equation}
where $P_{\text{eq}} = 0.791 \times 33 = 26.10$~fsw and
\[
  t_{1/2}^{\text{eff}} =
  \begin{cases}
    60 & \text{if } t_{1/2} \le 60\;\text{min}, \\
    t_{1/2} & \text{if } t_{1/2} > 60\;\text{min}.
  \end{cases}
\]

\textbf{Physical justification.}
This is a conservative modeling choice, not a physiological law.
Capping fast-compartment washout at the 60-minute rate prevents the model
from crediting rapid off-gassing in short surface intervals---a regime
where the Haldane exponential model is known to be less reliable and where
bubble dynamics (not captured by the model) may delay off-gassing.  The
rule was validated empirically in the Phase~I and Phase~II test series
(Hamilton et al., Chs.~IV--V).

\subsection{Repetitive dive NDL}

A repetitive dive is simply a standard NDL calculation
(Eq.~\ref{eq:ndl}) with modified initial loadings $P_0$ reflecting
the residual nitrogen from prior dives.  No special formula is needed;
the full compartment state carries all the information.


% ==================================================================
\section{Dive Profile Modeling}
% ==================================================================

\subsection{Complete dive profile}

A square dive profile consists of:
\begin{enumerate}[label=(\arabic*)]
  \item \textbf{Descent} at 60~fsw/min from surface to depth $d$,
        modeled with the Schreiner equation~\eqref{eq:schreiner}
        ($R = f_{\mathrm{N}_2} \times 60$).
  \item \textbf{Bottom time} at constant depth $d$, modeled with
        the Haldane equation~\eqref{eq:haldane}.
        Total bottom time is measured from start of descent to
        start of final ascent.
  \item \textbf{Safety stop} at 15~fsw (5~m) for 3~minutes (recommended
        on all dives; required near NDL or below 100~fsw).
  \item \textbf{Ascent} at 30~fsw/min (9~m/min) to the surface,
        modeled with the Schreiner equation ($R < 0$).
\end{enumerate}

\subsection{Bottom time definition}

``Bottom time is the total time in minutes from the beginning of descent
until the beginning of final ascent'' (PADI RDP instructions).  Descent
time is included; ascent time and safety-stop time are not.


% ==================================================================
\section{Published Table NDLs vs.\ Model Output}
% ==================================================================

When the model is evaluated at published table depths with a clean diver,
the computed NDLs agree with the published PADI RDP to within a few
minutes at most depths, but with a systematic pattern of residuals:

\begin{center}
\begin{tabular}{@{}rrrrr@{}}
\toprule
Depth & Published & Calculated & Diff & Controlling \\
(msw) & NDL (min) & NDL (min) & (min) & compartment \\
\midrule
10 & 219 & 262 & $+43$ & \#7 (80 min) \\
12 & 147 & 148 & $+1$  & \#6 (60 min) \\
14 &  98 &  99 & $+1$  & \#5 (40 min) \\
18 &  56 &  56 & $\phantom{+}0$ & \#4 (30 min) \\
22 &  37 &  37 & $\phantom{+}0$ & \#3 (20 min) \\
25 &  29 &  29 & $\phantom{+}0$ & \#3 (20 min) \\
30 &  20 &  19 & $-1$  & \#2 (10 min) \\
35 &  14 &  12 & $-2$  & \#1 (5 min)  \\
40 &   9 &   9 & $\phantom{+}0$ & \#1 (5 min)  \\
\bottomrule
\end{tabular}
\end{center}

The 10~msw outlier (published 219 vs.\ calculated 262) is a practical cap;
219 minutes ($\sim\!3.6$ hours) already exceeds realistic air supply.

The remaining residuals (conservative at shallow depths, slightly generous
at deep depths) are the signature of the empirical smoothing described in
Section~\ref{sec:m0derivation}.  The $M_0$ values were back-calculated
from a smooth NDL curve, so round-tripping (forward-computing NDLs from
the back-fitted $M_0$ values) does not reproduce the curve exactly because
the controlling compartment changes with depth.

\begin{remark}[Independent verification]
An independent implementer (``likedonuts'' on ScubaBoard, 2008) obtained
the same calculated NDLs to within rounding.  Dr.~Michael Powell confirmed
in the same thread that the discrepancies are due to curve smoothing by
Dr.~Rogers.
\end{remark}


% ==================================================================
\section{Safety Margins and Limitations}
% ==================================================================

\subsection{What the model does not capture}

The Haldane model is a dissolved-gas bookkeeping system, not a physiological
simulation.  It does not model:
\begin{itemize}[nosep]
  \item Bubble formation, growth, or transport
  \item Venous-to-arterial shunting (e.g., via patent foramen ovale)
  \item Tissue-specific blood flow variations with exercise or temperature
  \item Individual physiological variability
  \item Counter-diffusion effects (irrelevant for single-gas air diving)
\end{itemize}

\subsection{WXYZ rules}

For multi-day or intensive diving, the RDP imposes additional constraints:
if the ending pressure group after any dive is W or X, the minimum surface
interval before all subsequent dives is 1 hour; if Y or Z, the minimum
is 3 hours.  These rules protect slower compartments that the 60-minute
pressure-group system does not track directly.

\subsection{Conservatism}

The model embeds conservatism at multiple levels: Spencer's Doppler-based
NDL reductions below USN limits; the 60-minute surface washout cap; the
required safety stop; and the WXYZ rules.  Despite this, DCS can never be
completely eliminated---only made acceptably improbable (Hamilton et al.,
pp.~8--9).


% ==================================================================
\section{Summary of the Algorithm}
% ==================================================================

\begin{enumerate}[label=\textbf{Step \arabic*:}]
  \item \textbf{Initialize.}  Set all 14 compartment loadings to
    $P_0 = 0.791 \times 33 = 26.10$~fsw.

  \item \textbf{Descent.}  Apply Schreiner equation~\eqref{eq:schreiner}
    with $R = 0.791 \times 60 = 47.46$~fsw/min for time
    $t_d = d/60$ minutes.

  \item \textbf{Bottom.}  Apply Haldane equation~\eqref{eq:haldane}
    at constant $P_i = 0.791 \times (d + 33)$ for time
    $t_b = (\text{bottom time}) - t_d$.

  \item \textbf{Ascent / safety stop.}  Apply Schreiner equation with
    $R < 0$ at 30~fsw/min, including 3~min hold at 15~fsw.

  \item \textbf{NDL check.}  Before the dive, verify that the planned
    bottom time does not exceed the NDL (Eq.~\ref{eq:ndl}) for any
    compartment.

  \item \textbf{Pressure group.}  Read the 60-minute compartment loading
    and assign a letter A--Z.

  \item \textbf{Surface interval.}  Apply washout (Eq.~\ref{eq:washout})
    with the 60-min cap rule.

  \item \textbf{Repeat} from Step~2 for subsequent dives.
\end{enumerate}


% ==================================================================
\section*{References}
% ==================================================================

\begin{description}[style=nextline,leftmargin=0pt,font=\normalfont]

\item[Hamilton RW, Rogers RE, Powell MR, Vann RD (1994)]
  \emph{Development and validation of no-stop decompression procedures
  for recreational diving: The DSAT Recreational Dive Planner.}
  Santa Ana, CA: Diving Science and Technology Corp.

\item[Spencer MP (1976)]
  Decompression limits for compressed air determined by ultrasonically
  detected blood bubbles.
  \emph{J.\ Appl.\ Physiol.}\ 40(2):229--235.

\item[Workman RD (1965)]
  Calculation of decompression schedules for nitrogen-oxygen and
  helium-oxygen dives.  Research Report 6-65.
  Washington: Navy Experimental Diving Unit.

\item[Schreiner HR, Kelley PL (1971)]
  A pragmatic view of decompression.
  In: Lambertsen CJ, ed.\ \emph{Underwater Physiology IV.}
  New York: Academic Press.

\item[Hempleman HV (1975)]
  Decompression procedures for deep, open sea operations.
  In: \emph{Proc.\ 4th Symp.\ Underwater Physiology.}
  New York: Academic Press.

\item[Baker EC]
  Calculating no stop time for scuba diving.
  Available at: scubatechphilippines.com.

\item[Baker EC]
  Understanding $M$-values.
  Available at: scubatechphilippines.com.

\end{description}

\end{document}